% Standalone problem document, generated from the adjacent README.md.
% Compile with XeLaTeX. Mathematical content is maintained in Markdown.
\documentclass[11pt,a4paper]{article}
\usepackage[fontsize=10.5pt]{scrextend}
\usepackage[margin=22mm,headheight=15pt,headsep=9mm,footskip=11mm]{geometry}
\usepackage{fontspec}
\setmainfont{texgyrepagella}[Extension=.otf,UprightFont=*-regular,BoldFont=*-bold,ItalicFont=*-italic,BoldItalicFont=*-bolditalic]
\setsansfont{texgyreheros}[Extension=.otf,UprightFont=*-regular,BoldFont=*-bold,ItalicFont=*-italic,BoldItalicFont=*-bolditalic]
\setmonofont{texgyrecursor}[Extension=.otf,UprightFont=*-regular,BoldFont=*-bold,ItalicFont=*-italic,BoldItalicFont=*-bolditalic,Scale=MatchLowercase]
\usepackage{amsmath,amssymb}
\usepackage{unicode-math}
\setmathfont{texgyrepagella-math.otf}
\usepackage{xcolor,microtype,enumitem,needspace,fancyhdr}
\usepackage{bookmark}
\definecolor{ink}{HTML}{17344B}
\definecolor{accent}{HTML}{197D80}
\definecolor{muted}{HTML}{596773}
\hypersetup{colorlinks=true,linkcolor=accent,urlcolor=accent,pdftitle={MF-24 - A
dimension-independent polynomial norm bound from super-identical
pseudospectra},pdfauthor={Open Problems in Numerical Linear Algebra}}
\urlstyle{same}
\setlength{\parindent}{0pt}
\setlength{\parskip}{5pt plus 1pt minus 1pt}
\linespread{1.04}
\setlength{\emergencystretch}{3em}
\widowpenalty=10000
\clubpenalty=10000
\displaywidowpenalty=10000
\allowdisplaybreaks[1]
\setlist{nosep,leftmargin=1.5em}
\providecommand{\tightlist}{\setlength{\itemsep}{0pt}\setlength{\parskip}{0pt}}
\providecommand{\pandocbounded}[1]{#1}
\pagestyle{fancy}
\fancyhf{}
\fancyhead[L]{\sffamily\footnotesize\color{muted}OPEN PROBLEMS IN NUMERICAL LINEAR ALGEBRA}
\fancyhead[R]{\sffamily\bfseries\footnotesize\color{accent}MF-24}
\fancyfoot[L]{\parbox[b]{0.9\textwidth}{\sffamily\fontsize{8}{10}\selectfont\color{muted}Editorial ratings; a bounded search cannot certify that a problem remains unsolved.\\Literature check: 2026-09-15}}
\fancyfoot[R]{\sffamily\footnotesize\color{muted}\thepage}
\renewcommand{\headrulewidth}{0.3pt}
\renewcommand{\headrule}{\hbox to\headwidth{\color{accent}\leaders\hrule height \headrulewidth\hfill}}
\makeatletter
\renewcommand{\section}{\Needspace{5\baselineskip}\@startsection{section}{1}{0pt}{12pt plus 2pt minus 2pt}{4pt}{\sffamily\bfseries\large\color{ink}}}
\renewcommand{\subsection}{\Needspace{4\baselineskip}\@startsection{subsection}{2}{0pt}{9pt plus 1pt minus 1pt}{3pt}{\sffamily\bfseries\normalsize\color{ink}}}
\makeatother
\setcounter{secnumdepth}{0}
\begin{document}
{\sffamily\small\color{muted}Matrix functions and stability\par}
\vspace{3pt}
{\raggedright\hyphenpenalty=10000\exhyphenpenalty=10000\sffamily\bfseries\fontsize{21}{24}\selectfont\color{ink}A
dimension-independent polynomial norm bound from super-identical
pseudospectra\par}
\vspace{7pt}
{\sffamily\small\textbf{Difficulty:} challenging\quad\textbf{Importance:} interesting
to the community\par}
{\sffamily\small\bfseries\color{accent}Status: Lean verified\par}
\vspace{4pt}
{\color{accent}\hrule height 0.8pt}
\vspace{6pt}
\textbf{Topic:} Polynomial matrix norms and pseudospectral information

\textbf{Rating rationale:} Existing information controls polynomial
norms with a dimension-dependent constant; a uniform comparison would
strengthen pseudospectral inference about nonnormal dynamics and
polynomial matrix functions.

\subsection{Resolution: negative, 15 September
2026}\label{resolution-negative-15-september-2026}

\textbf{Solved negatively by Georg Maierhofer (University of
Cambridge).} The manuscript is dated 14 September 2026 and passed a
fresh independent informal AI-agent audit on 15 September 2026. The
exact dimension-dependent constants and their optimal growth remain
open.

The
\href{https://github.com/ajt60gaibb/OpenProblemsInNLA/blob/main/references/mf24-counterexample/proof.pdf}{counterexample
manuscript}, \textbf{Theorem 1 and equation (15)}, constructs real
nonnegative nilpotent weighted shifts of order \(N=(m+1)^2\), for every
integer \(m\ge2\) and real \(t>1\). All singular values agree after
every complex scalar shift. For the common polynomial
\(p_m(z)=\sum_{j=1}^m z^{(m+2)j}\), it proves

\[\frac{\|p_m(X_{m,t})\|_2}{\|p_m(Y_{m,t})\|_2}
\ge \frac{\sqrt m}{1+(m-1)/t^2},\qquad p_m(Y_{m,t})\ne0.\]

Taking \(t=m\) gives a ratio at least \((4/5)\sqrt m\), so the original
uniform-boundedness target is refuted. \textbf{Corollary 3} gives
\(C_N\ge\sqrt{\lfloor\sqrt N\rfloor-1}\) for \(N\ge9\). It does not
determine the sharp \(C_N\) or close the gap to the known upper bound.

\textbf{Evidence level:} a complete written argument with reproducible
exact and numerical checks, developed and self-reviewed with AI
assistance. The
\href{https://github.com/ajt60gaibb/OpenProblemsInNLA/blob/main/reviews/2026-09-15-pr264/README.md}{fresh
maintainer audit} confirms the complete negative resolution and records
independent mathematical review and fresh checks. The
\href{https://github.com/ajt60gaibb/OpenProblemsInNLA/blob/main/references/mf24-counterexample/independent-review.md}{submitted
informal Codex AI-agent audit} is also retained. Those original audits
were informal; the separate Lean verification is documented below. No
external human peer review is asserted. See the
\href{https://github.com/ajt60gaibb/OpenProblemsInNLA/blob/main/references/mf24-counterexample/README.md}{source,
checks and provenance}. The existing difficulty and importance ratings
are retained as historical ratings of the original boundedness question.

\newpage

\subsection{Lean proof and verification
evidence}\label{lean-proof-and-verification-evidence}

\textbf{Formalization: George Stepaniants, Department of Computing and
Mathematical Sciences, California Institute of Technology.} Georg
Maierhofer retains credit for the mathematical counterexample above. The
\href{https://github.com/sgstepaniants/OpenProblemsInNLA/blob/208e30d80f73ef661b1f019c7de93c70254f7e48/matrix-functions-and-stability/MF-24/lean/Solution.lean}{complete
Lean proof at immutable revision 208e30d8} pins Lean \textbf{4.33.1},
Mathlib \textit{0df444a360eaa60ab8c11dca51a86af692955474} and LeanCert
\textit{621a43d7cf21f87872392a01e874f2f1dbddc926}.

The declarations in namespace \textit{NLA.MF24} include:

\begin{itemize}
\tightlist
\item
  \textit{family\_super\_identical}: equality of every ordered singular
  value after every complex shift.
\item
  \textit{polynomial\_norm\_ratio} and
  \textit{finite\_rational\_family\_ratio}: the explicit norm-ratio
  estimates for the original weighted-shift family.
\item
  \textit{arbitrarily\_large\_ratios}: a finite-dimensional pair and
  polynomial exceeding every proposed real comparison bound.
\item
  \textit{no\_uniform\_comparison}: the complete negative answer to the
  original uniform-boundedness question.
\end{itemize}

All \textbf{22 declarations}, including their full determinant, spectrum
and operator-norm dependencies, are listed in
\href{https://github.com/ajt60gaibb/OpenProblemsInNLA/blob/main/matrix-functions-and-stability/MF-24/lean/formalization.yaml}{formalization.yaml}
and checked against the independently reviewed
\href{https://github.com/ajt60gaibb/OpenProblemsInNLA/blob/main/matrix-functions-and-stability/MF-24/lean/Challenge.lean}{Challenge}.
The
\href{https://github.com/ajt60gaibb/OpenProblemsInNLA/blob/main/matrix-functions-and-stability/MF-24/lean/NLA/MF24/Definitions.lean}{definitions}
use genuine complex polynomial evaluation and the induced norm on
complex Euclidean space. Singular-value equality includes multiplicities
and zeros; no normality, fixed dimension or restricted shift is assumed.
The formal estimate uses the conservative denominator \(t^2+m\), giving
\((2/3)\sqrt m\) at \(t=m\). This still refutes the complete original
target. The manuscript's sharper constant, dimension-padding corollary
and optimal growth question are not claimed as Lean-verified results.

On \textbf{15 September 2026},
\href{https://github.com/sgstepaniants/OpenProblemsInNLA/actions/runs/35033148310/job/104596164346}{the
actual GitHub Linux verification job} checked this exact proof commit.
All 22 exports passed LeanCert kernel assertions, real non-root
Comparator statement matching, default-kernel replay and transitive
axiom checks. Only \textit{propext}, \textit{Classical.choice} and
\textit{Quot.sound} support the target results. Invalid-proof,
mismatched-statement, forbidden-axiom and sandbox controls passed. The
\href{https://github.com/ajt60gaibb/OpenProblemsInNLA/blob/main/matrix-functions-and-stability/MF-24/lean/verification/linux-2026-09-15}{retained
raw logs, artifact and input hashes} bind that execution to all 192
candidate project files. These are actual GitHub Linux checks; no local
macOS Lean execution is claimed.

Two independent AI-agent referees accepted the complete mathematical
source and reviewed the actual runtime evidence:
\href{https://github.com/ajt60gaibb/OpenProblemsInNLA/blob/main/matrix-functions-and-stability/MF-24/lean/reviews/final/referee-root/REVIEW.md}{referee
1} and
\href{https://github.com/ajt60gaibb/OpenProblemsInNLA/blob/main/matrix-functions-and-stability/MF-24/lean/reviews/final/referee-inequalities/REVIEW.md}{referee
2}. They applied the repository's scoped Tau Ceti protocol. This is
separate from external human peer review or certification that the
checker software is infallible.

To reproduce the ordinary proof build, enter
\href{https://github.com/ajt60gaibb/OpenProblemsInNLA/blob/main/matrix-functions-and-stability/MF-24/lean/README.md}{the
Lean project directory} with Lean installed and run:

\begin{verbatim}
lake exe cache get && lake build
\end{verbatim}

The default target is the complete Solution. For the additional
sandboxed Comparator, kernel replay and rejection checks, follow the
\href{https://github.com/ajt60gaibb/OpenProblemsInNLA/blob/main/docs/lean/README.md}{shared
Linux verification instructions}. Solution never imports the Challenge's
deliberate specification placeholders. This evidence names the exact
tested proof revision; subsequent documentation commits have separate
published-commit checks.

\newpage

\subsection{Context and notation}\label{context-and-notation}

All matrices are complex and all norms are spectral norms. Write
\(s_1(M)\geq\cdots\geq s_N(M)\) for the singular values of \(M\). Define

\[A\sim_{\rm sip}B\quad\Longleftrightarrow\quad
s_j(A-zI)=s_j(B-zI)\quad(z\in\mathbb C,\ 1\leq j\leq N).\]

Define the sharp comparison constant

\[C_N=\sup\left\{
\frac{\|p(A)\|_2}{\|p(B)\|_2}:
A,B\in\mathbb C^{N\times N},\ A\sim_{\rm sip}B,
p\in\mathbb C[z],\ p(B)\ne0
\right\}.\]

As recorded in
\href{https://github.com/ajt60gaibb/OpenProblemsInNLA/blob/main/eigenvalues-and-inverse-problems/SP-15/README.md}{SP-15},
super-identical pseudospectral data also imply ordinary similarity, so
\(p(A)=0\) if and only if \(p(B)=0\).

\subsection{Problem statement}\label{problem-statement}

Determine whether \(\sup_{N\geq1}C_N<\infty\). Equivalently, does there
exist an absolute constant \(C\) such that

\[\|p(A)\|_2\leq C\|p(B)\|_2\]

for every dimension, every super-identical-pseudospectral pair and every
polynomial? Determining the sharp \(C_N\) is the closely related
quantitative question, retained here rather than split into another
entry.

\subsection{References}\label{references}

\begin{itemize}
\item
  Fortier Bourque and Ransford,
  \href{https://doi.org/10.1112/jlms/jdn085}{\emph{Super-identical
  pseudospectra}}, \textbf{p.~513, immediately after Theorem 1.3},
  explicitly ask both whether their \(\sqrt N\) bound is optimal and
  whether a dimension-independent bound exists. The sharp-constant
  notation above is editorial; the boundedness question is
  source-stated.
\item
  Thomas Ransford and Nathan Walsh,
  \href{https://arxiv.org/pdf/2109.14472v2}{\emph{A four-mean theorem
  and its application to pseudospectra}, arXiv:2109.14472v2}, \textbf{9
  July 2022}, \textbf{Theorem 1.3}, prove

  \[\|p(A)\|_2<\sqrt{N-2}\,\|p(B)\|_2\qquad(N\geq4)\]

  unless both sides vanish. \textbf{Proposition 5.1} and its following
  argument show sharpness for \(N=4\): \(C_4=\sqrt2\). \textbf{Theorem
  1.4} concerns unbounded condition numbers of similarity transforms in
  fixed dimension, a different quantity.
\end{itemize}

Thus \(C_N\leq\sqrt{N-2}\) for \(N\geq4\). The small-dimensional
unitary/transpose classification gives \(C_1=C_2=C_3=1\).

\subsection{Original scope and literature
check}\label{original-scope-and-literature-check}

The following is the repository's pre-claim literature record, retained
unchanged.

The unknown absolute constant, not the superseded \(\sqrt N\) sharpness
guess, is the admission target. This comparison does not follow just
from ordinary similarity, since a badly conditioned similarity can
change norms substantially. Searches through 2026-09-11 for
super-identical polynomial norm bounds, dimension-independent bounds,
sharp constants and successors to the four-mean paper found no full
solution. The current arXiv version history still ends at v2, 9 July
2022. This is an older-source question supported by bounded searches; no
2026 explicit open-status reaffirmation was located.
\end{document}
